\subsubsection{一个无解的偏微分方程}\label{subsec:pde-without-solutions}

不时会出现标题类似"求解所有偏微分方程的一种方法"的论文. 这类论文的内容总是远远达不到标题中的宣称. 据说, 一篇这样的论文启发 Lewy 构造了他著名的完全无解的线性偏微分方程例子. 这个例子也突出了 Cauchy-Kovalevskaya 结果中解析性的重要性.

\begin{Theorem}
\label{thm:lewy-nonsolvability}
对复值函数 $u(x,y,z)$, 令
\begin{equation}\label{eq:lewy-differential-operator}
Lu=-u_x-iu_y+2i(x+iy)u_z.
\end{equation}
那么存在实值函数 $f(x,y,z)$, 属于 $C^\infty(\mathbb{R}^3)$ 类, 使方程
\begin{equation}\label{eq:lewy-inhomogeneous-equation}
Lu=f(x,y,z)
\end{equation}
在任意开子集 $\Omega\subset\mathbb{R}^3$ 内都没有 $C^1(\Omega)$ 类的解.
\end{Theorem}

注意, 当 $f$ 解析时, 可以应用 Cauchy-Kovalevskaya 定理, \eqref{eq:lewy-inhomogeneous-equation} 的非特征初值问题具有局部解. 相比之下, 对非解析的 $f$, 即使不指定初始条件, 也可能无解.

不在这里给出定理的完整证明, 而只概述其中一些主要思想. 首先证明以下引理.

\begin{Lemma}
\label{lem:lewy-local-nonsolvability}
设 $\psi\in C^\infty(\mathbb{R})$ 为实值函数, 且 $\psi$ 在 $z_0$ 处不实解析. 那么方程
\begin{equation}\label{eq:lewy-local-obstruction-equation}
Lu=\psi'(z)
\end{equation}
在点 $(0,0,z_0)$ 的任意邻域内都没有 $C^1$ 类的解.
\end{Lemma}

\begin{Proof}
反设存在解 $u$, 定义在 $(0,0,z_0)$ 的一个邻域内, 例如在 $x^2+y^2<\epsilon$, $|z-z_0|<\epsilon$ 时有定义. 令
\begin{equation}\label{eq:lewy-polar-transformed-function}
v(r,\theta,z)=e^{i\theta}\sqrt r\,u(\sqrt r\cos\theta,\sqrt r\sin\theta,z).
\end{equation}
经一些代数运算可得, $v$ 满足方程
\begin{equation}\label{eq:lewy-polar-differential-equation}
-2v_r-\dfrac irv_\theta+2iv_z=\psi'(z).
\end{equation}
令
\begin{equation}\label{eq:lewy-angular-average}
V(z,r)=\int_0^{2\pi}v(r,\theta,z)\,\mathrm{d}\theta,
\end{equation}
并将 \eqref{eq:lewy-polar-differential-equation} 对 $\theta$ 积分, 得到
\begin{equation}\label{eq:lewy-averaged-differential-equation}
V_z+iV_r=-\pi i\psi'(z).
\end{equation}
显然, $V$ 在 $0<r<\epsilon$, $|z-z_0|<\epsilon$ 时属于 $C^1$ 类, 并在 $0\leq r<\epsilon$, $|z-z_0|<\epsilon$ 时连续. 此外, $V(z,0)=0$. 现在考虑 $W=V(z,r)+\pi i\psi(z)$. 由 \eqref{eq:lewy-averaged-differential-equation} 得到 $W_z+iW_r=0$, 从而 $W$ 是区域 $|z-z_0|<\epsilon$, $0<r<\epsilon$ 内复变量 $z+ir$ 的全纯函数. 此外, $W$ 连续到边界 $r=0$, 且在那里为纯虚值. 由 Schwarz 反射原理 (参见 Problem~\ref{ex:schwarz-reflection-principle}), 定义 $W(z,-r)=-\overline{W(z,r)}$, 即可将 $W$ 延拓成 $|r|<\epsilon$, $|z-z_0|<\epsilon$ 上的全纯函数. 但这意味着 $W(z,0)$ 的虚部 $\pi\psi(z)$ 实解析.
\end{Proof}

方程
\begin{equation}\label{eq:lewy-shifted-forcing}
Lu=\psi'(z-2y_0x+2x_0y)
\end{equation}
通过简单代换
\begin{equation}\label{eq:lewy-coordinate-translation}
U(x,y,z)=u(x+x_0,y+y_0,z+2y_0x-2x_0y).
\end{equation}
化为 \eqref{eq:lewy-local-obstruction-equation}. 因而, 给定任意点 $(x_0,y_0,z_0)$, 都能找到函数 $f(x,y,z)$, 使 \eqref{eq:lewy-inhomogeneous-equation} 在 $(x_0,y_0,z_0)$ 的任意邻域内无解.

现在考虑有限个点 $(x_i,y_i,z_i)$, $i=1,\ldots,k$. 将构造函数 $f$, 使 \eqref{eq:lewy-inhomogeneous-equation} 在这些点中任意一点的邻域内都无解. 为此, 在 $C^\infty(\mathbb{R})$ 中选取处处不解析的实值函数 $\psi$ (参见 Problem~\ref{ex:smooth-nowhere-analytic-function}), 并作设定
\begin{equation}\label{eq:lewy-finitely-many-forcing-centres}
f(x,y,z)=\sum_{i=1}^kc_i\psi'(z-2y_ix+2x_iy),
\end{equation}
其中 $c_i$ 为实系数. 那么, 对 $(c_2,c_3,\ldots,c_k)$ 的任意选择, 至多存在一个值 $c_1=c_1(c_2,\ldots,c_k)$, 使 \eqref{eq:lewy-inhomogeneous-equation} 在 $(x_1,y_1,z_1)$ 的某个邻域内有解. 反设有两个这样的值 $c_1$ 和 $\tilde c_1$. 则可推出方程
\begin{equation}\label{eq:lewy-two-coefficient-contradiction}
Lu=(c_1-\tilde c_1)\psi'(z-2y_1x+2x_1y)
\end{equation}
在 $(x_1,y_1,z_1)$ 的一个邻域内有解, 与上面证明的结果矛盾.

类似地, 给定 $(c_1,c_3,\ldots,c_k)$ 时, 至多有一个值 $c_2=c_2(c_1,c_3,\ldots,c_k)$, 使 \eqref{eq:lewy-inhomogeneous-equation} 在 $(x_2,y_2,z_2)$ 的某个邻域内可能有解. 现在将 $c_i$ 限制在含有 $l$ 个元素的集合中, 其中 $l>k$. $(c_1,\ldots,c_k)$ 有 $l^k$ 种可能选择. 然而, 使关系 $c_1=c_1(c_2,\ldots,c_k)$, $c_2=c_2(c_1,c_3,\ldots,c_k)$ 等中的任意一个成立的选择至多有 $kl^{k-1}$ 种. 因而存在这样的选择, 使 \eqref{eq:lewy-inhomogeneous-equation} 在各点 $(x_i,y_i,z_i)$ 的任意邻域内都无解. 事实上, 如果 $l$ 相对于 $k$ 很大, 则 $c_i$ 的"大多数"选择都具有这一性质.

要完成定理的证明, 需将这一论证推广到可数个点. 设 $(x_i,y_i,z_i)$, $i\in\mathbb{N}$, 为在 $\mathbb{R}^3$ 中稠密的点列. 然后作设定
\begin{equation}\label{eq:lewy-countably-many-forcing-centres}
f(x,y,z)=\sum_{i=1}^{\infty}c_i\psi'(z-2y_ix+2x_iy).
\end{equation}
如果 $c_i$ 在 $i\to\infty$ 时充分迅速地趋于零, 那么 $f$ 属于 $C^\infty$ 类. 仍可证明, 对 $c_i$ 的"大多数"选择, \eqref{eq:lewy-inhomogeneous-equation} 在任意点 $(x_i,y_i,z_i)$ 的任意邻域内都无解; 因而它在任何地方都无解. 完成这一论证并不容易, 需要泛函分析方法. 这里不再追究这一点, 而请读者参阅文献, 例如 \cite{Jo}.

\paragraph*{Problems}

\begin{enumerate}[label=\thesection.\arabic*.,ref=\thesection.\arabic*,start=7,series=book-problems]

\item\label{ex:multi-index-geometric-series-identity} 若对 $i=1,\ldots,n$ 有 $|x_i|<1$, 证明
\begin{equation}\label{eq:multivariable-geometric-product}
\sum_\alpha\bm{x}^{\alpha}=\dfrac1{(1-x_1)(1-x_2)\cdots(1-x_n)}.
\end{equation}
对两端施加 $D^\beta$ 并比较. 然后令 $\bm{x}=(q,q,\ldots,q)$, 由此将得到 \eqref{eq:multi-index-geometric-derivative-sum}.

\item\label{ex:complete-power-series-facts} 补全上文关于幂级数的三个断言的证明中的其余细节 (与单变量幂级数的情形类似地进行).

\item\label{ex:taylor-series-real-analyticity} 设 Taylor 级数
\begin{equation}\label{eq:exercise-local-taylor-analyticity}
\sum_\alpha c_\alpha\bm{x}^{\alpha}
\end{equation}
在 $D=\{\bm{x}\in\mathbb{R}^n\mid|\bm{x}|<r\}$ 中收敛. 证明该 Taylor 级数表示的函数在 $D$ 中实解析.

\item\label{ex:analytic-taylor-convex-domain} 设 $f$ 在包含原点的凸区域 $D$ 内实解析. 令
\begin{equation}\label{eq:exercise-convex-domain-taylor-series}
\sum_\alpha c_\alpha\bm{x}^{\alpha}
\end{equation}
为相应的 Taylor 级数. 假设此级数在某个 $\bm{x}\in D$ 处绝对收敛. 证明它收敛到 $f(\bm{x})$. Hint: 对 $0<\lambda<1$ 考虑 $f(\lambda\bm{x})$, 再取极限 $\lambda\to1$.

\item\label{ex:elliptic-cauchy-instability} 求初值问题
\begin{equation}\label{eq:elliptic-cauchy-high-frequency-data}
u_{xx}+u_{yy}=0,\quad u(x,0)=0,\quad u_y(x,0)=\dfrac{\sin nx}{n}.
\end{equation}
的解. 讨论 $n\to\infty$ 时会发生什么. 与 $u_{xx}-u_{yy}=0$ 的类似初值问题比较.

\item\label{ex:nonlinear-analytic-breakdown} 考虑初值问题
\begin{equation}\label{eq:nonlinear-cauchy-characteristic-problem}
u_y=uu_x,\qquad u(x,0)=f(x),
\end{equation}
其中 $f$ 解析. 证明, 除非 $f$ 单调递减, 否则不可能对所有正的 $y$ 都存在解析 (甚至连续) 解. Hint: $u$ 沿由 $\mathrm{d}x/\mathrm{d}y=-u$ 给出的特征为常数. 如果 $f$ 不单调递减, 特征必相交.

\item\label{ex:reduced-initial-problem-equivalence} 详细证明, \eqref{eq:first-order-scalar-nonlinear-system} 配以初始条件 \eqref{eq:first-order-scalar-nonlinear-initial-data}, 等价于 \eqref{eq:differentiated-scalar-nonlinear-equation} 配以初始条件 \eqref{eq:differentiated-scalar-initial-data}. 单独的 \eqref{eq:first-order-scalar-nonlinear-system} 是否等价于 \eqref{eq:differentiated-scalar-nonlinear-equation}? 如果不等价, 它们有何区别?

\item\label{ex:laplace-system-standard-form} 方程组 \eqref{eq:first-order-laplace-system} 不具有 \eqref{eq:ck-intro-quasilinear-system} 的形式. 讨论如何将 \eqref{eq:first-order-laplace-system} 的初值问题化为标准形式 \eqref{eq:ck-intro-quasilinear-system}, \eqref{eq:ck-intro-zero-initial-data}.

\item\label{ex:schwarz-reflection-principle} 设 $\Omega$ 为 $\mathbb{C}$ 中关于实轴对称的区域. 设 $f$ 在 $\Omega\cap\{\operatorname{Im}z>0\}$ 内全纯, 在 $\Omega\cap\{\operatorname{Im}z\geq0\}$ 上连续. 此外, 假设 $f$ 在 $\Omega\cap\{\operatorname{Im}z=0\}$ 上取实值. 证明, 对 $\operatorname{Im}z<0$ 定义 $f(z)=\overline{f(\bar z)}$, 就可将 $f$ 延拓为整个 $\Omega$ 内的全纯函数. Hint: 证明, 对任意闭的可求长曲线 $C$, 若 $C$ 及其内部均位于 $\Omega$ 内, 则 $\int_Cf(z)\,\mathrm{d}z=0$. 只需对 $C$ 为三角形的情形证明这一点.

\item\label{ex:surface-descriptions-equivalent} 证明 $C^k$ (或解析) 曲面的三个定义确实等价.

\item\label{ex:smooth-nowhere-analytic-function} 证明函数
\begin{equation}\label{eq:smooth-nowhere-analytic-series}
f(x)=\sum_{n=1}^{\infty}\dfrac{\cos(n!x)}{(n!)^n}
\end{equation}
在 $\mathbb{R}$ 上属于 $C^\infty$ 类, 但处处不实解析. Hint: 先证明对任意 $M$ 和 $r$, $f$ 都不属于 $C_{M,r}(0)$. 再证明对任意有理数 $q$, $f(x)-f(x+2\pi q)$ 都解析.

\end{enumerate}
